\documentclass[conference]{IEEEtran}
\usepackage[utf8]{inputenc}
\usepackage[T1]{fontenc}
\usepackage{amsmath}
\usepackage{amssymb}
\usepackage{graphicx}
\usepackage{url}

\begin{document}

\title{Implémentation du jeu Hex via la Programmation Orientée Objet et stratégies de décision en C++}

\author{
    \IEEEauthorblockN{1\textsuperscript{st} Jair Anderson Vasquez Torres}
    \IEEEauthorblockA{\textit{Ingénieur Degree Programme STIC} \\
    \textit{ENSTA Paris}\\
    Paris, France \\
    jair-anderson.vasquez@ensta.fr}
    \and
    \IEEEauthorblockN{2\textsuperscript{nd} Santiago Florido Gomez}
    \IEEEauthorblockA{\textit{Ingénieur Degree Programme STIC} \\
    \textit{ENSTA Paris}\\
    Paris, France \\
    santiago.florido@ensta-paris.fr}
}

\maketitle

\begin{abstract}
Ce projet implémente en C++ un système complet pour le jeu Hex, conçu principalement comme un exercice de Programmation Orientée Objet. La solution organise le domaine du jeu au moyen de classes qui encapsulent le plateau, l’état, les coordonnées et les règles, permettant la génération de coups légaux et la vérification de victoire au moyen d’un parcours BFS sur les connexions des pions. Sur cette base s’intègrent des stratégies de décision (p. ex., Negamax avec hachage et table de transposition) et une évaluation de positions découplée du moteur, qui peut être heuristique ou s’appuyer sur un modèle neuronal de valeur intégré à l’exécutable (exportable en TorchScript), sans lier la conception à une architecture spécifique. De plus, une interface graphique (GUI) en SFML a été développée pour faciliter l’interaction, la visualisation du plateau et les tests du comportement des agents.
\end{abstract}

\begin{IEEEkeywords}
Hex game, C++, Object-Oriented Programming, Negamax, Transposition Table, TorchScript, SFML.
\end{IEEEkeywords}



\section{Principes de la classification bayésienne}

Les méthodes d’apprentissage bayésien partent d’un ensemble de données d’observation noté \(X\) et d’un ensemble de classes \(c\). Dans ce cadre, une observation est modélisée sous forme vectorielle, et l’on peut définir une probabilité a priori d’observer le vecteur \(X\), notée \(P(X)\), ainsi qu’une probabilité a priori d’appartenance à une classe \(c\), notée \(P(c)\), comme indiqué dans \eqref{eq:priorX} et \eqref{eq:priorc}. On peut également exprimer la probabilité conjointe \(P(X,c)\), donnée en \eqref{eq:joint}.

\begin{equation}
\label{eq:priorX}
P(X)
\end{equation}

\begin{equation}
\label{eq:priorc}
P(c)
\end{equation}

\begin{equation}
\label{eq:joint}
P(X,c)
\end{equation}

En utilisant le concept de probabilité conditionnelle, il est possible d’exprimer la probabilité d’observer les données \(X\) sachant que l’observation appartient à la classe \(c\). Cette probabilité conditionnelle, notée \(P(X\mid c)\), est aussi appelée vraisemblance (conditionnelle), comme en \eqref{eq:likelihood}.

\begin{equation}
\label{eq:likelihood}
P(X\mid c)
\end{equation}

La probabilité conjointe peut alors s’écrire comme le produit de la vraisemblance et de l’a priori de classe, ce qui donne \eqref{eq:joint_factorization}.

\begin{equation}
\label{eq:joint_factorization}
P(X,c)=P(X\mid c)\,P(c)
\end{equation}

Enfin, en appliquant le théorème de Bayes, on obtient une probabilité a posteriori \(P(c\mid X)\), qui exprime la probabilité qu’une observation \(X\) appartienne à une classe spécifique \(c\) dans le domaine du système. Cette relation est donnée par \eqref{eq:bayes}.

\begin{equation}
\label{eq:bayes}
P(c\mid X)=\frac{P(X\mid c)\,P(c)}{P(X)}
\end{equation}


L’apprentissage bayésien consiste donc en la détermination empirique des lois de probabilité, référée à partir d’un ensemble de données statistiques, réalisée sur un ensemble de données d’apprentissage étiquetées \(\{X_i,\;c(X_i)\}\). La détermination de la classe \(c\) donne lieu à deux formes de sélection en fonction du critère d’optimisation employé.

\subsection{Maximum A Posteriori (MAP)}

On se focalise sur un critère de maximisation de la probabilité a posteriori, de sorte qu’à partir d’une observation \(X\), on cherche la sélection de la classe \(c\) la plus probable \cite{onera_ensta_estimateur_bayesien_slides}, suivant l’expression \eqref{eq:cmap}.

\begin{equation}
\label{eq:cmap}
c_{\mathrm{MAP}}^{\ast}(X)=\arg\max_{c\in C} P(c\mid X).
\end{equation}

En utilisant des logarithmes (utile numériquement et pour l’analyse), le critère MAP peut s’écrire comme indiqué en \eqref{eq:cmap_log}.

\begin{equation}
\label{eq:cmap_log}
c^{*}_{\mathrm{MAP}}(X)=\arg\max_{c\in C}\Bigl(\log P(X\mid c)+\log P(c)\Bigr).
\end{equation}

La probabilité a priori peut être estimée à partir de la base de données d’entraînement par la relation \eqref{eq:prior_est}, où \(N_c\) désigne le nombre d’exemples appartenant à la classe \(c\), défini en \eqref{eq:Nc_def}.

\begin{equation}
\label{eq:prior_est}
P(c)=\frac{N_c}{N}.
\end{equation}

\begin{equation}
\label{eq:Nc_def}
N_c=\left|\left\{\,i:\;c_i=c\,\right\}\right|.
\end{equation}

De son côté, pour la vraisemblance, on ajuste les données observées à un modèle \(P(X\mid c)\) pour chaque classe (par exemple : histogramme, Naive Bayes, ou modèle gaussien).


\subsection{Maximum Vraisemblance (ML)}
De son côté, le critère de ML (Maximum Likelihood) cherche la sélection de la classe qui rend la donnée \(X\) d’intérêt pour la classification la plus vraisemblable, au moyen de l’expression de maximisation \eqref{eq:cml}.

\begin{equation}
\label{eq:cml}
c_{\mathrm{ML}}^{\ast}(X)=\arg\max_{c\in C} P(X\mid c).
\end{equation}

Il suit ainsi un principe de recherche de la classe pour laquelle la donnée observée \(X\) \cite{yuille_bayes_decision_theory}.

Dans la représentation logarithmique de la maximisation, il devient évident que l’on conserve l’effet de la vraisemblance, et qu’à la différence du MAP il n’y a pas d’influence de la distribution a priori de la classe, ce qui conduit à \eqref{eq:cml_log}.

\begin{equation}
\label{eq:cml_log}
c_{\mathrm{ML}}^{\ast}(X)=\arg\max_{c\in C}\log P(X\mid c).
\end{equation}

Cela rend compte du fait que la principale différence entre ML et MAP est le déplacement des limites de sélection à cause de la nature de la fonction a priori pour chaque classe. Ainsi, dans le cas d’une classification à deux classes, ML utiliserait le seuil \(0\) défini en \eqref{eq:ml_threshold}.

\begin{equation}
\label{eq:ml_threshold}
\log P(X\mid c_1)-\log P(X\mid c_2) > 0.
\end{equation}

Et pour MAP, on obtient le critère \eqref{eq:map_threshold}.

\begin{equation}
\label{eq:map_threshold}
\log P(X\mid c_1)-\log P(X\mid c_2) > \log\frac{P(c_2)}{P(c_1)}.
\end{equation}

Cela implique que si \(P(c_1)\) est grande, alors \(\log\frac{P(c_2)}{P(c_1)}\) est négatif, et MAP a besoin de moins d’évidence (en termes de likelihood) pour choisir \(c_1\). On conclut donc que MAP, à la différence de ML, considère dans la classification le fait que certaines classes, par leur propre dynamique et nature, sont plus probables que d’autres, ce qui peut produire un effet en classification qui n’est pas désiré lorsque les classes sont, par exemple, déséquilibrées.



\begin{thebibliography}{00}
    
\bibitem{onera_ensta_estimateur_bayesien_slides}
ONERA and ENSTA ParisTech, ``Estimation et Identification statistiques --- Cours 2 : Estimateur Bayésien,'' course slides (PDF), Sept. 13, 2023, accessed Feb. 18, 2026. [Course material]. Provided by the user.

\bibitem{yuille_bayes_decision_theory}
A. Yuille, ``Lecture 2. Bayes Decision Theory,'' course Stat 161/261, Johns Hopkins University, Spring 2014, accessed Feb. 18, 2026. [Online]. Available: \url{https://www.cs.jhu.edu/~ayuille1/courses/Stat161-261-Spring14/RevisedLectureNotes2.pdf}

\end{thebibliography}

\end{document}
